\subsection{Hilbertian tensor product of hilbertian spaces}

DEFINITION 1. --- Let \(E_{1}, \ldots, E_{n}\) be hilbertian spaces. The completion of the Hausdorff prehilbertian space \(E_{1} \otimes_{2} \ldots \otimes_{2} E_{n}\) is called the hilbertian tensor product of the \(E_{i}\) and is denoted by \(E_{1} \hat{\otimes}_{2} \ldots \hat{\otimes}_{2} E_{n}\) (or \(\hat{\otimes}_{2} E_{i}\) ).

Let \(F_{1}, \ldots, F_{n}\) be hilbertian spaces and \(u_{i} \in \mathcal{L}(E_{i}, F_{i})\) for \(1 \leqslant i \leqslant n\) . The continuous linear mapping \(u_{1} \otimes \ldots \otimes u_{n}\) then extends to a continuous linear mapping \(u_{1} \hat{\otimes}_{2} \ldots \hat{\otimes}_{2} u_{n}\) from \(E_{1} \hat{\otimes}_{2} \ldots \hat{\otimes}_{2} E_{n}\) into \(F_{1} \hat{\otimes}_{2} \ldots \hat{\otimes}_{2} F_{n}\) . We have

\[
\| u _ {1} \hat {\otimes} _ {2} \dots \hat {\otimes} _ {2} u _ {n} \| = \| u _ {1} \| \dots \| u _ {n} \|\tag{10}
\]

by formula (8) of V, p.~27. Moreover, if \(1_{E}\) denotes the identity mapping of any hilbertian space E, we have

\[
1 _ {\mathrm{E} _ {1}} \hat {\otimes} _ {2} \dots \hat {\otimes} _ {2} 1 _ {\mathrm{E} _ {n}} = 1 _ {\mathrm{E}} \quad \text { with } \quad \mathrm{E} = \mathrm{E} _ {1} \hat {\otimes} _ {2} \dots \hat {\otimes} _ {2} \mathrm{E} _ {n}.\tag{11}
\]

Finally, if \(\mathbf{G}_1, \ldots, \mathbf{G}_n\) are hilbertian spaces and \(v_i \in \mathcal{L}(\mathbf{F}_i; \mathbf{G}_i)\) for \(1 \leqslant i \leqslant n\) , we get

\[
\left(v _ {1} \circ u _ {1}\right) \hat {\otimes} _ {2} \dots \hat {\otimes} _ {2} \left(v _ {n} \circ u _ {n}\right) = \left(v _ {1} \hat {\otimes} _ {2} \dots \hat {\otimes} _ {2} v _ {n}\right) \circ \left(u _ {1} \hat {\otimes} _ {2} \dots \hat {\otimes} _ {2} u _ {n}\right).\tag{12}
\]

We leave to the reader the task of formulating the «commutativity» and the «associativity» of the hilbertian tensor product, in analogy with what has been said above for prehilbertian spaces.

Remark. --- Let \(E_{1}, \ldots, E_{n}\) be Hausdorff prehilbertian spaces, and \(\hat{E}_{1}, \ldots, \hat{E}_{n}\) their respective completions. Then \(E_{1} \otimes_{2} \ldots \otimes_{2} E_{n}\) is a prehilbertian subspace of \(\hat{E}_{1} \otimes_{2} \ldots \otimes_{2} \hat{E}_{n}\) . Since the mapping \((x_{1}, \ldots, x_{n}) \mapsto x_{1} \otimes \ldots \otimes x_{n}\) from \(\hat{E}_{1} \times \cdots \times \hat{E}_{n}\) into \(\hat{E}_{1} \otimes_{2} \ldots \otimes_{2} \hat{E}_{n}\) is continuous, \(E_{1} \otimes_{2} \ldots \otimes_{2} E_{n}\) is dense in \(\hat{E}_{1} \otimes_{2} \ldots \otimes_{2} \hat{E}_{n}\) . A fortiori the completion of \(E_{1} \otimes_{2} \ldots \otimes_{2} E_{n}\) is precisely the hilbertian space \(\hat{E}_{1} \otimes_{2} \ldots \otimes_{2} \hat{E}_{n}\) . This completion is sometimes simply written as \(E_{1} \otimes_{2} \ldots \otimes_{2} E_{n}\) (or \(\bigotimes_{1 \leq i \leq n} E_{i}\) ).

PROPOSITION 3. --- Let \(E_{1}, \ldots, E_{n}\) be hilbertian spaces. Suppose that for \(1 \leqslant i \leqslant n\) the space \(E_{i}\) is a hilbertian sum of a family \((\mathrm{E}_{i,\alpha})_{\alpha \in \mathbf{A}(i)}\) of closed vector subspaces. Then \(E_{1} \otimes_{2} \ldots \otimes_{2} E_{n}\) is a hilbertian sum of the family of subspaces \(E_{1,\alpha_{1}} \hat{\otimes}_{2} \ldots \hat{\otimes}_{2} E_{n,\alpha_{n}}\) with \((\alpha_{1}, \ldots, \alpha_{n})\) ranging over \(A(1) \times \cdots \times A(n)\) .

By formula (6) of V, p.~27, the subspaces \(\mathrm{E}_{1,\alpha_1}\hat{\otimes}_2\ldots \hat{\otimes}_2\mathrm{E}_{n,\alpha_n}\) of \(\mathrm{E}_1\hat{\otimes}_2\ldots \hat{\otimes}_2\mathrm{E}_n\) are mutually orthogonal. For every integer \(i\) between 1 and \(n\) , the set \(\bigcup_{\alpha \in \mathbf{A}(i)}\mathrm{E}_{i,\alpha}\) is total in \(\mathrm{E}_i\) , and the multilinear mapping \((x_{1},\dots,x_{n})\mapsto x_{1}\otimes \dots \otimes x_{n}\) is continuous. It follows that the union of the subspaces \(\mathrm{E}_{1,\alpha_1}\hat{\otimes}_2\ldots \hat{\otimes}_2\mathrm{E}_{n,\alpha_n}\) is total, hence prop. 3.

COROLLARY 1. --- For \(1 \leqslant i \leqslant n\) , let \((e_{i,\alpha})_{\alpha \in \mathbf{A}(i)}\) be an orthonormal basis of \(\mathbf{E}_i\) . Then the family of vectors \(e_{1,\alpha_1} \otimes \ldots \otimes e_{n,\alpha_n}\) as \((\alpha_1, \ldots, \alpha_n)\) ranges over \(\mathbf{A}(1) \times \cdots \times \mathbf{A}(n)\) is an orthonormal basis of \(\mathbf{E}_1 \hat{\otimes}_2 \ldots \hat{\otimes}_2 \mathbf{E}_n\) .

COROLLARY 2. --- Let \(E_{1}\) and \(E_{2}\) be two hilbertian spaces, and \((e_{i})_{i\in I}\) an orthonormal basis of \(E_{1}\) . Let \((y_{i})_{i\in I}\) be a family of elements of \(E_{2}\) , such that \(\sum_{i\in I}\|y_{i}\|^{2}<+\infty\) . Then the family \((e_{i}\otimes y_{i})_{i\in I}\) is summable in \(E_{1}\hat{\otimes}_{2}E_{2}\) ; moreover, every element of \(E_{1}\hat{\otimes}_{2}E_{2}\) can be written uniquely in the form \(\sum_{i\in I}e_{i}\otimes y_{i}\) with \(\sum_{i\in I}\|y_{i}\|^{2}<+\infty\) .

Let \(F_{i}\) be the line in \(E_{1}\) generated by the \(e_{i} (i \in I)\) . Then \(E_{1}\) is the hilbertian sum of the family of subspaces \((F_{i})_{i \in I}\) . By prop. 3, the space \(E_{1} \hat{\otimes} E_{2}\) is the hilbertian sum of the family of subspaces \((F_{i} \hat{\otimes}_{2} E_{2})_{i \in I}\) , hence cor. 2 follows.

Examples. --- 1) By cor. 1, the space \(\ell^2 (\mathrm{I})\hat{\otimes}_2\ell^2 (\mathrm{J})\) is canonically isomorphic to \(\ell^2 (\mathrm{I}\times \mathrm{J})\) , the tensor product \(\mathbf{x}\otimes \mathbf{y}\) of \(\mathbf{x} = (x_i)_{i\in \mathrm{I}}\) and \(\mathbf{y} = (y_j)_{j\in \mathrm{J}}\) can be identified with the family \((x_iy_j)_{i\in \mathrm{I},j\in \mathrm{J}}\) . Similarly, by cor. 2, \(\ell^2 (\mathrm{I})\hat{\otimes}_2\mathrm{E}\) can be identified with \(\ell_{\mathrm{E}}^{2}(\mathrm{I})\) , in such a way that we have \((x_{i})_{i\in \mathrm{I}}\otimes y = (x_{i}y)_{i\in \mathrm{I}}\) for every \(\mathbf{y}\) in the hilbertian space E.

* 2) Let X be a Hausdorff topological space, and \(\mu\) a positive measure on X. Let E be a hilbertian space. We can identify \(L^2(X, \mu) \otimes_2 E\) with \(L_E^2(X, \mu)\) in a canonical way: if \(f\) is the class of the square integrable scalar function \(f\) on X, and if \(a\) belongs to E, then \(f \otimes a\) is the class of the function \(x \mapsto f(x).a\) with values in E.

Let Y be a Hausdorff topological space and v a positive measure on Y. In an analogous manner, we can identify the hilbertian spaces \(\mathrm{L}^{2}(\mathrm{X},\mu)\hat{\otimes}_{2}\mathrm{L}^{2}(\mathrm{Y},v)\) and \(\mathrm{L}^{2}(\mathrm{X}\times\mathrm{Y},\mu\otimes v)\) ; then \(\dot{f}\otimes\dot{g}\) can be identified with the class of the function \((x,y)\mapsto f(x)g(y)\) on \(X\times Y\) .

